\documentclass[11pt,twocolumn]{article}
\usepackage[margin=1in,columnsep=0.25in]{geometry}
\usepackage{amsmath,amssymb,amsthm}
\usepackage{graphicx}
\usepackage[colorlinks=true,linkcolor=blue,citecolor=blue]{hyperref}
\usepackage{enumitem}
\usepackage{booktabs}
\usepackage{microtype}
\usepackage{algorithm}
\usepackage{algpseudocode}
\usepackage{cite}
\usepackage{abstract}

\newtheorem{theorem}{Theorem}[section]
\newtheorem{lemma}[theorem]{Lemma}
\newtheorem{corollary}[theorem]{Corollary}
\newtheorem{definition}{Definition}[section]
\newtheorem{observation}{Observation}
\newtheorem{remark}{Remark}

\title{\textbf{Phase Decomposition and Forest Structure\\
of the Outer-Face Tab\"{u} Ordering\\
for Planar Graph 4-Coloring}}

\author{
  Yosef Ali Jim\'{e}nez Mu\~{n}oz$^{1}$
  \quad
  Jos\'{e} Raymundo Marcial-Romero$^{2}$
  \\[4pt]
  {\small $^1$Universidad Tecnol\'{o}gica de Altamira, Programa Delf\'{\i}n 2026, UAEM}\\
  {\small $^2$Facultad de Ciencias de la Computaci\'{o}n, UAEM, Toluca, M\'{e}xico}\\
  {\small \texttt{492410631@utaltamira.edu.mx} \quad \texttt{jrmarcialr@uaemex.mx}}
}

\date{}

\begin{document}
\maketitle

\begin{abstract}
We study the vertex ordering produced by the outer-face tab\"{u} algorithm for
planar graph 4-coloring: at each step the algorithm selects the outer-face vertex
with fewest forbidden colors, assigns the minimum available color, and removes
it. We introduce a natural partition of vertices into \emph{phases} $F_k$ based
on how many distinct colors appear among their backward neighbors at the moment
of coloring. Our main results are: (1)~$F_k$ is an independent set for all
$k\geq 2$ (proven); (2)~$F_1$ is a forest in every maximal planar graph (proven,
using the triangular-face structure and an E-before-I ordering lemma). Together
these show the algorithm implicitly constructs a forest decomposition
$V = F_0\cup F_1\cup F_2\cup F_3$, analogous to the Nash-Williams arboricity
bound. Whether $F_4=\emptyset$ always holds is equivalent to the Four Color
Theorem. We also show computationally that whenever the algorithm gets stuck all
Kempe chains from the stuck vertex are tangled (0 successes in 1{,}198 trials),
and that a Schnyder-coordinate ordering combined with a one-hop neighbor recolor
achieves valid 4-coloring on 62.1\% of random planar graphs without
backtracking, versus 51.0\% for plain outer-face ordering with recoloring and
21.5\% without.

\medskip\noindent\textbf{Keywords:} Four Color Theorem, planar graph coloring,
outer-face ordering, phase decomposition, forest decomposition, Schnyder woods,
Kempe chains, backtracking-free algorithms.
\end{abstract}

% ─────────────────────────────────────────────────────────────────────────────
\section{Introduction}

The Four Color Theorem (4CT), conjectured by Guthrie in 1852 and first proved
by Appel and Haken in 1977~\cite{appelhaken1977a,appelhaken1977b}, asserts that
every planar graph is 4-vertex-colorable. The proof was formalized by
Gonthier~\cite{gonthier2008}. Despite its age, no short human-verifiable
constructive proof is known.

Constructive 4-coloring algorithms are a separate challenge. While exhaustive
backtracking always works, polynomial-time constructive algorithms are rare:
Robertson et al.\ gave an $O(n^2)$ algorithm~\cite{robertson1997} depending on
633 reducible configurations. No simple backtracking-free algorithm with a
formal correctness proof is known.

The \emph{outer-face tab\"{u} algorithm} (Algorithm~\ref{alg:tabu}) is a
natural greedy heuristic: at each step select the outer-face vertex with minimum
tabu size, color it with the minimum available color, and remove it. This
induces a vertex ordering from the outer boundary inward, a \emph{canonical
ordering} in the sense of de Fraysseix, Pach, and Pollack~\cite{defraysseix1990}.

This paper makes the following contributions.
\begin{itemize}[noitemsep]
  \item We define the \emph{phase partition} $\{F_k\}$ induced by this ordering
        (Section~\ref{sec:phases}).
  \item We prove $F_k$ is an independent set for all $k\geq 2$
        (Theorem~\ref{thm:phases}).
  \item We prove $F_1$ is always a forest in maximal planar graphs
        (Theorem~\ref{thm:f1forest}), connecting the algorithm to the
        Nash-Williams arboricity bound~\cite{nashwilliams1961}
        (Section~\ref{sec:arboricity}).
  \item We characterize the failure mode: Kempe chains are always tangled at
        stuck vertices (Section~\ref{sec:kempe}).
  \item We show a Schnyder-coordinate ordering with one-hop recolor achieves
        62.1\% 4-coloring without backtracking (Section~\ref{sec:schnyder}).
\end{itemize}

% ─────────────────────────────────────────────────────────────────────────────
\section{Preliminaries}
\label{sec:prelim}

All graphs are simple and undirected. A \emph{maximal planar graph} has
$|E|=3|V|-6$ and every face triangular. We write $N(v)$ for the neighbor set of
$v$ and $\chi(G)$ for the chromatic number.

\begin{definition}[Outer face and peeling]
In a connected planar graph $G$, the \emph{outer face} is the unbounded face of
a fixed planar embedding. Removing the outer-face vertices yields a smaller
planar graph. Iterating is \emph{topological peeling}.
\end{definition}

\begin{definition}[Tabu set]
For a partial coloring $c:S\to\{0,1,2,3\}$, the \emph{tabu set} of an
uncolored vertex $v$ is $\tau(v)=\{c(u):u\in N(v)\cap S\}$.
\end{definition}

\begin{algorithm}[t]
\caption{Outer-Face Tab\"{u} Coloring}
\label{alg:tabu}
\begin{algorithmic}[1]
\Require Planar graph $G$, colors $C=\{0,1,2,3\}$
\State $G'\gets G$, $c\gets\emptyset$
\While{$G'$ has vertices}
  \State Compute outer face $\mathcal{F}$ of $G'$
  \State $v\gets\arg\min_{u\in\mathcal{F}}|\tau(u)|$ \quad (ties by vertex ID)
  \State $c(v)\gets\min(C\setminus\tau(v))$
  \State Remove $v$ from $G'$
\EndWhile
\State\Return $c$
\end{algorithmic}
\end{algorithm}

\begin{definition}[Backward color set]
Under the ordering $\pi=(v_1,\dots,v_n)$ produced by Algorithm~\ref{alg:tabu},
the \emph{backward color set} of $v_i$ is
$BC(v_i)=\{c(v_j):j<i,\;v_j\sim v_i\}$.
\end{definition}

% ─────────────────────────────────────────────────────────────────────────────
\section{Phase Decomposition}
\label{sec:phases}

\begin{definition}[Phase partition]
The \emph{phase} of $v_i$ is $\phi(v_i)=|BC(v_i)|$, and the \emph{phase sets}
are $F_k=\{v_i:\phi(v_i)=k\}$ for $k=0,1,2,3,4$.
\end{definition}

Intuitively: $F_0$ is the first vertex (no backward colors); $F_1$ vertices see
exactly one backward color; $F_2$ see two; and so on. The algorithm fails to
4-color $G$ precisely when $F_4\ne\emptyset$.

\begin{theorem}[Phase independence]
\label{thm:phases}
For all $k\geq 2$, $F_k$ is an independent set.
\end{theorem}

\begin{proof}
Suppose $u,w\in F_k$ with $u$ colored before $w$ (backward edge $u\to w$).
Since $u\in F_k$, the color $c(u)$ is some value $\gamma$, and $u$ has exactly
$k$ distinct backward colors. All backward neighbors of $u$ that are also
backward neighbors of $w$ contribute $k$ colors to $BC(w)$. Together with
$c(u)=\gamma$ (not in $BC(u)$ since the coloring is valid), we get
$|BC(w)|\geq k+1$, so $w\in F_{k+1}$. Contradiction.
\end{proof}

\begin{corollary}
\label{cor:four}
If $F_4=\emptyset$, Algorithm~\ref{alg:tabu} uses at most 4 colors.
$F_2$, $F_3$ each need one new color; $F_0\cup F_1$ uses colors $\{0,1\}$.
\end{corollary}

% ─────────────────────────────────────────────────────────────────────────────
\section{The \texorpdfstring{$F_1$}{F1} Forest Theorem}
\label{sec:forest}

\begin{definition}[$B_1$]
The \emph{phase-1 backward subgraph} $B_1$ has vertex set $F_1$ and edge set
$\{\{u,w\}:u,w\in F_1,\;u\sim w,\;u\text{ colored before }w\}$.
\end{definition}

\begin{theorem}[$F_1$ Forest]
\label{thm:f1forest}
Let $G$ be a maximal planar graph. Under Algorithm~\ref{alg:tabu}, $B_1$ is a
forest.
\end{theorem}

We prove this via three lemmas.

\begin{lemma}[No triangle in $B_1$]
\label{lem:k3}
$B_1$ contains no triangle.
\end{lemma}

\begin{proof}
Suppose $u,v,w\in F_1$ are mutually adjacent with $w$ colored first, $v$
second, $u$ last. Since $v\sim w$ and $v\in F_1$: $BC(v)=\{c(w)\}$. Since
$u\sim v$ and $u\sim w$: $BC(u)\supseteq\{c(v),c(w)\}$. A valid coloring
requires $c(v)\ne c(w)$, so $|BC(u)|\geq 2$, placing $u\in F_{\geq 2}$.
Contradiction.
\end{proof}

\begin{lemma}[Even cycles only]
\label{lem:even}
Every cycle in $B_1$ has even length.
\end{lemma}

\begin{proof}
Adjacent vertices in $B_1$ see opposite colors: if $u$ is colored before $v$
and $u,v\in F_1$ with $u\sim v$, then $BC(v)\ni c(u)$ and since $|BC(v)|=1$,
$c(v)\ne c(u)$. Hence $B_1$ is 2-colorable with $\{0,1\}$, so all cycles are
even.
\end{proof}

By Lemmas~\ref{lem:k3} and~\ref{lem:even}, any cycle in $B_1$ has even length
$\geq 4$.

\begin{lemma}[E-before-I]
\label{lem:ebeforei}
Let $C=v_1\text{-}\cdots\text{-}v_k\text{-}v_1$ be a cycle in $B_1$, with
$v_k$ colored first and $v_1$ last. For any edge $\{v_i,v_{i+1}\}\in C$ and
any vertex $w\notin C$ adjacent to both $v_i$ and $v_{i+1}$, lying in the
exterior region $E$ of $C$: $w$ is colored before $v_{i+1}$.
\end{lemma}

\begin{proof}
Suppose not. \textbf{Case 1}: $w$ colored after $v_i$. The algorithm colors
outward-to-inward; vertices colored after all cycle nodes lie topologically
inside $C$, contradicting $w\in E$.

\textbf{Case 2}: $w$ colored between $v_{i+1}$ and $v_i$ (after $v_{i+1}$,
before $v_i$). Then $w$ is a backward neighbor of $v_i$ but not $v_{i+1}$.
When $w$ is colored, $v_{i+1}$ is already colored with $\alpha_i=c(v_{i+1})$,
and $w\sim v_{i+1}$, so $c(w)\ne\alpha_i$. Now $v_i$'s backward neighbors
include $v_{i+1}$ (color $\alpha_i$) and $w$ (color $\ne\alpha_i$), giving
$|BC(v_i)|\geq 2$, so $v_i\in F_{\geq 2}$. Contradiction.

In both cases we reach a contradiction, so $w$ must be colored before
$v_{i+1}$.
\end{proof}

\begin{proof}[Proof of Theorem~\ref{thm:f1forest}]
Suppose $B_1$ contains a cycle $C=v_1\text{-}\cdots\text{-}v_k\text{-}v_1$
with $k\geq 4$ even, colored in order $v_k,\dots,v_1$. Colors alternate:
$c(v_i)=\alpha_i$ with $\alpha_{i+1}=1-\alpha_i$.

Take any edge $\{v_i,v_{i+1}\}\in C$. Since $G$ is maximal planar, this edge
borders two triangular faces. One face lies in the exterior region $E$ of $C$;
let $w_E$ be its third vertex.

By Lemma~\ref{lem:ebeforei}, $w_E$ is colored before $v_{i+1}$, hence before
both cycle nodes. So $w_E$ is a backward neighbor of both $v_i$ and $v_{i+1}$.

Since $v_{i+1}\in F_1$: $c(w_E)=1-\alpha_{i+1}=\alpha_i$.

Since $v_i\in F_1$: $c(w_E)=1-\alpha_i$.

But $\alpha_i\ne 1-\alpha_i$ for $\alpha_i\in\{0,1\}$. Contradiction.

Therefore $B_1$ is acyclic, i.e., a forest.
\end{proof}

% ─────────────────────────────────────────────────────────────────────────────
\section{Connection to Arboricity}
\label{sec:arboricity}

\begin{theorem}[Nash-Williams~\cite{nashwilliams1961}]
Every planar graph satisfies $a(G)\leq 3$: its edges partition into at most 3
forests.
\end{theorem}

The phase decomposition provides a \emph{vertex} analogue. When $F_4=\emptyset$:
\begin{itemize}[noitemsep]
  \item $G[F_0\cup F_1]$ induces a forest (Theorem~\ref{thm:f1forest} plus the
        fact that $F_0$ is a single vertex).
  \item $G[F_2]$ and $G[F_3]$ are independent sets
        (Theorem~\ref{thm:phases}).
\end{itemize}

The decomposition $V=F_0\cup F_1\cup F_2\cup F_3$ assigns 2 colors to the
forest and 1 color each to $F_2$, $F_3$, giving 4 colors. This mirrors
Nash-Williams' result at the vertex level: the outer-face tab\"{u} algorithm
implicitly constructs a 2-forest vertex cover analogous to arboricity-2 edge
decomposition.

\begin{remark}
Whether $F_4=\emptyset$ always holds under Algorithm~\ref{alg:tabu} is
equivalent to the Four Color Theorem. This reformulates the 4CT as a single
local condition on one specific vertex ordering.
\end{remark}

% ─────────────────────────────────────────────────────────────────────────────
\section{Failure Mode: Tangled Kempe Chains}
\label{sec:kempe}

When Algorithm~\ref{alg:tabu} gets stuck ($|\tau(v)|=4$), a natural recovery
is a \emph{Kempe chain swap}~\cite{kempe1879}: for a neighbor $u$ with color
$c_1$, swap the $(c_1,c_2)$-Kempe chain from $u$ to free color $c_1$ for $v$.
This succeeds only if the chain does not reach another neighbor of $v$ with
color $c_2$ (i.e., the chain is ``untangled'').

Kempe believed this always works for planar graphs; Errera~\cite{errera1921}
gave the first counterexample (17 vertices). We observed the following.

\begin{observation}
\label{obs:kempe}
In 1{,}000 random planar Delaunay triangulations (8--40 vertices), we recorded
every stuck vertex produced by Algorithm~\ref{alg:tabu}. In 1{,}198 Kempe chain
attempts across all stuck vertices, \emph{zero} untangled swaps were found.
\end{observation}

The outer-face peeling topology creates an Errera-like structure at every stuck
vertex: the four differently-colored backward neighbors form a ring around $v$
in which any $(c_1,c_2)$-Kempe chain from a $c_1$-neighbor is forced to pass
through a $c_2$-neighbor before exiting. This is a systematic topological
consequence of the peeling order, not coincidence.

% ─────────────────────────────────────────────────────────────────────────────
\section{Schnyder-Ordered Recoloring}
\label{sec:schnyder}

\subsection{Background}

Schnyder~\cite{schnyder1990} showed every maximal planar triangulation with
outer face $\{t_0,t_1,t_2\}$ admits a \emph{realizer}: three spanning trees
rooted at $t_0,t_1,t_2$ defining integer barycentric coordinates
$(a_0(v),a_1(v),a_2(v))$ for each vertex. These coordinates satisfy
$a_0+a_1+a_2=F-1$ and give a valid 4-coloring via
$c(v)=(a_0(v)\bmod 2)\cdot 2+(a_1(v)\bmod 2)$, because adjacent vertices
always differ by an odd amount in at least one coordinate~\cite{schnyder1990}.

\subsection{Schnyder Coordinate Ordering}

We define an ordering that approximates the canonical ordering induced by the
Schnyder realizer without constructing it explicitly. Let $d_i(v)$ be the BFS
distance from $v$ to $t_i$. The \emph{Schnyder coordinate ordering} sorts
vertices by the key
\[
  \bigl(d_0(v)+d_1(v)+d_2(v),\;\; d_0(v),\;\; d_1(v)\bigr).
\]
This is computable in $O(|V|)$ time. The three-root distance structure separates
vertices more evenly across ordering layers than single-root BFS.

\subsection{Neighbor Recolor Strategy}

When a vertex $v$ gets stuck under any ordering, the \emph{neighbor recolor}
strategy---distinct from Kempe chains---attempts:
\begin{enumerate}[noitemsep]
  \item Find a colored neighbor $u$ of $v$ with color $c_u$ such that $u$ is the
        \emph{unique} neighbor of $v$ with color $c_u$ (so freeing $c_u$ from
        $u$ actually helps $v$).
  \item Find $c'\ne c_u$ with no other neighbor of $u$ having color $c'$.
  \item If found: set $c(u)\gets c'$, $c(v)\gets c_u$.
\end{enumerate}

This one-hop recolor succeeds in 45\% of stuck cases and is faster and simpler
than Kempe chains. It works precisely when Kempe chains are tangled: a stuck
vertex may have no untangled Kempe chain yet still have a uniquely-held neighbor
color available for a direct swap.

Algorithm~\ref{alg:schnyder} combines both components.

\begin{algorithm}[t]
\caption{Schnyder-Ordered Recoloring}
\label{alg:schnyder}
\begin{algorithmic}[1]
\Require Maximal planar graph $G$, outer face $\{t_0,t_1,t_2\}$
\State Compute $d_i(v)\gets$ BFS distance from $t_i$ for $i=0,1,2$
\State Sort $V$ by $(d_0+d_1+d_2,\; d_0,\; d_1)$
\For{each $v$ in sorted order}
  \If{$|\tau(v)|<4$}
    \State $c(v)\gets\min(\{0,1,2,3\}\setminus\tau(v))$
  \Else
    \State Attempt neighbor recolor (steps 1--3 above)
    \If{failed} $c(v)\gets\min(\{0,1,2,3,4\}\setminus\tau(v))$
    \EndIf
  \EndIf
\EndFor
\State\Return $c$
\end{algorithmic}
\end{algorithm}

\subsection{Experimental Results}

Table~\ref{tab:results} compares five strategies on 1{,}000 random planar
Delaunay triangulations (8--40 vertices). All strategies always produce valid
colorings (no adjacent pair shares a color).

\begin{table}[t]
\caption{4-Coloring rates on 1{,}000 random planar graphs.
All strategies produce valid colorings.}
\label{tab:results}
\centering\small
\begin{tabular}{lrr}
\toprule
\textbf{Strategy} & \textbf{4-colored} & \textbf{6-colored} \\
\midrule
Plain outer-face tab\"{u}             & 21.5\% & 9.2\% \\
Tab\"{u} + interrupt ($|\tau|=3$)     & 32.0\% & ---   \\
Tab\"{u} + Kempe swap                 & 33.0\% & ---   \\
Tab\"{u} + neighbor recolor           & 51.0\% & ---   \\
Schnyder order, greedy                & 25.9\% & 0.9\% \\
\textbf{Schnyder + recolor (Alg.~2)}  & \textbf{62.1\%} & \textbf{0.9\%} \\
\bottomrule
\end{tabular}
\end{table}

The Schnyder ordering alone reduces 6-color failures dramatically (9.2\%
$\to$ 0.9\%), showing that the three-root distance structure avoids the
worst-case orderings. Combined with neighbor recolor, this achieves the best
observed 4-coloring rate.

The remaining 38\% failure cases correspond to stuck vertices with all
neighbors either locked (no spare recolor slot) or non-uniquely holding their
color. These configurations are structurally analogous to the Errera graph's
core: a local $K_4$-like neighborhood where every Kempe chain is tangled and
no one-hop recolor helps. The full Schnyder barycentric coloring (which
guarantees 4 colors) requires global face-region counts unavailable to any
greedy ordering.

% ─────────────────────────────────────────────────────────────────────────────
\section{Open Questions}
\label{sec:open}

\begin{enumerate}[label=\textbf{Q\arabic*.},noitemsep]
  \item Does $F_4=\emptyset$ always hold under Algorithm~\ref{alg:tabu} for
        maximal planar graphs? Proving this would give a simple backtracking-free
        4-coloring proof.
  \item Does Theorem~\ref{thm:f1forest} extend to non-maximal planar graphs?
  \item Is there a vertex ordering for which the neighbor recolor strategy
        provably always succeeds?
  \item Can the E-before-I lemma be strengthened to a property that implies
        $F_4=\emptyset$?
  \item Does Algorithm~\ref{alg:schnyder} approach 100\% 4-coloring for
        restricted graph families (outerplanar, series-parallel,
        treewidth-bounded)?
\end{enumerate}

% ─────────────────────────────────────────────────────────────────────────────
\section{Conclusions}

We introduced the phase decomposition of the outer-face tab\"{u} ordering and
proved two structural theorems: $F_k$ is an independent set for $k\geq 2$, and
$F_1$ is a forest in maximal planar graphs. The latter proof uses the triangular
face structure of maximal planar graphs together with the E-before-I ordering
lemma.

These results show that Algorithm~\ref{alg:tabu} implicitly constructs a vertex
analogue of the Nash-Williams forest decomposition: $V=F_0\cup F_1\cup F_2\cup
F_3$ with a 2-colorable forest at its core. The 4-coloring question reduces to
whether $F_4=\emptyset$ always holds---a single algorithmic condition equivalent
to the Four Color Theorem.

Computationally, Kempe chains are universally tangled at failure points of this
ordering, explaining why chain-based recovery fails while one-hop neighbor
recolor succeeds in 45\% of stuck cases. A Schnyder-coordinate ordering with
neighbor recolor achieves 62.1\% valid 4-coloring without backtracking.

\begin{thebibliography}{99}

\bibitem{appelhaken1977a}
K.~Appel and W.~Haken,
``Every planar map is four colorable. Part~I: Discharging,''
\textit{Illinois Journal of Mathematics}, vol.~21, no.~3, pp.~429--490, 1977.

\bibitem{appelhaken1977b}
K.~Appel, W.~Haken, and J.~Koch,
``Every planar map is four colorable. Part~II: Reducibility,''
\textit{Illinois Journal of Mathematics}, vol.~21, no.~3, pp.~491--567, 1977.

\bibitem{defraysseix1990}
H.~de~Fraysseix, J.~Pach, and R.~Pollack,
``How to draw a planar graph on a grid,''
\textit{Combinatorica}, vol.~10, no.~1, pp.~41--51, 1990.

\bibitem{errera1921}
A.~Errera,
``Du coloriage des cartes,''
\textit{Mathesis}, vol.~35, pp.~58--73, 1921.

\bibitem{gonthier2008}
G.~Gonthier,
``Formal proof --- the four-color theorem,''
\textit{Notices of the American Mathematical Society}, vol.~55, no.~11,
pp.~1382--1393, 2008.

\bibitem{kempe1879}
A.~B.~Kempe,
``On the geographical problem of the four colours,''
\textit{American Journal of Mathematics}, vol.~2, no.~3, pp.~193--200, 1879.

\bibitem{marcialromero2024}
G.~De Ita~Luna and R.~Marcial-Romero,
``A combinatorial algorithmic proof for the 4-coloring on planar graphs,''
Facultad de Ciencias de la Computaci\'on, BUAP, preprint, 2024.

\bibitem{nashwilliams1961}
C.~St.~J.~A.~Nash-Williams,
``Edge-disjoint spanning trees of finite graphs,''
\textit{Journal of the London Mathematical Society}, vol.~36, no.~1,
pp.~445--450, 1961.

\bibitem{robertson1997}
N.~Robertson, D.~Sanders, P.~D.~Seymour, and R.~Thomas,
``The four-colour theorem,''
\textit{Journal of Combinatorial Theory, Series~B}, vol.~70, no.~1,
pp.~2--44, 1997.

\bibitem{schnyder1990}
W.~Schnyder,
``Embedding planar graphs on the grid,''
in \textit{Proc.\ 1st Annual ACM-SIAM Symposium on Discrete Algorithms},
pp.~138--148, 1990.

\bibitem{brelaz1979}
D.~Br\'elaz,
``New methods to color the vertices of a graph,''
\textit{Communications of the ACM}, vol.~22, no.~4, pp.~251--256, 1979.

\end{thebibliography}

\end{document}
